%% ============================================================
%%  Asynchronous Post-Quantum Signature Verification for
%%  Low-Latency TLS 1.3 Handshakes
%%  Computer Networks — Elsevier
%%  Revised Manuscript
%% ============================================================
\documentclass[preprint,10pt,3p]{elsarticle}

%% ---- Core packages -----------------------------------------
\usepackage[T1]{fontenc}
\usepackage[utf8]{inputenc}
\usepackage{amsmath,amssymb}
\usepackage{graphicx}
\usepackage{booktabs}
\usepackage{multirow}
\usepackage{array}
\usepackage{float}
\usepackage{xcolor}
\usepackage[hidelinks]{hyperref}
\usepackage{url}
\usepackage[protrusion=true,expansion=false]{microtype}
\usepackage{enumitem}
\usepackage[ruled,vlined,linesnumbered]{algorithm2e}
\usepackage{placeins}   %% provides \FloatBarrier
\usepackage{geometry}
\geometry{margin=1.8cm, top=2cm, bottom=2cm, includehead=false}
\usepackage{svg} % in preamble
\usepackage{hyperref}


% Add these packages to your LaTeX preamble for SVG support with math expressions:

\usepackage{svg}           % For \includesvg command
\usepackage{amsmath}       % For math symbols
\usepackage{amssymb}       % For additional math symbols
\usepackage[utf8]{inputenc} % For Unicode support (if needed)
\usepackage[T1]{fontenc}   % For font encoding

% Optional: Set the SVG path (adjust to your directory structure)
% \svgpath{{./figures/}}

%% ---- Float placement tuning --------------------------------
\setcounter{topnumber}{9}
\setcounter{bottomnumber}{9}
\setcounter{totalnumber}{20}
\renewcommand{\topfraction}{0.90}
\renewcommand{\bottomfraction}{0.90}
\renewcommand{\textfraction}{0.05}
\renewcommand{\floatpagefraction}{0.85}

%% ---- Tight float/caption spacing ---------------------------
\setlength{\textfloatsep}{8pt plus 2pt minus 2pt}
\setlength{\floatsep}{6pt plus 2pt minus 2pt}
\setlength{\intextsep}{6pt plus 2pt minus 2pt}
\setlength{\abovecaptionskip}{4pt}
\setlength{\belowcaptionskip}{2pt}

%% ---- Paragraph spacing -------------------------------------
\setlength{\parskip}{2pt plus 1pt}

%% ---- Figure path -------------------------------------------
%% Place all matplotlib-generated figures inside ./figures/
\graphicspath{{figures/}}

%% ---- Convenience macros ------------------------------------
\newcommand{\Dt}{\ensuremath{\Delta t}}
\newcommand{\ms}{\,ms}

\journal{Computer Networks}

%% ============================================================
\begin{document}
\begin{frontmatter}

\title{Deferred Asynchronous Post-Quantum Signature Verification for
Low-Latency TLS~1.3 Handshakes}

% \author[au1]{Tahera Begum Abdul\corref{cor1}}
% %\author[au1]{ Tahera Begum Abdul \\ \texttt{taherabegum.abdul@gmail.com}}
% \ead{taherabegum.abdul@gmail.com}
% \author[au1]{K.~Venkata Ramana}\\

% %\href{mailto:taherabegum.abdul@gmail.com}{taherabegum.abdul@gmail.com}

% \address[au1]{Department of Computer Science and Systems Engineering,
% Andhra University, Visakhapatnam, Andhra Pradesh, India}

\author[au1]{Tahera Begum Abdul\corref{cor1}}
\ead{taherabegum.abdul@gmail.com}

\author[au1]{K.~Venkata Ramana}

\cortext[cor1]{Corresponding author}

\address[au1]{Department of Computer Science and Systems Engineering,
Andhra University, Visakhapatnam, Andhra Pradesh, India}


%% ---- Abstract ----------------------------------------------
\begin{abstract}
Post-quantum cryptography (PQC) standardisation by NIST has prompted urgent
migration efforts across Internet security infrastructure, with Transport Layer
Security (TLS) 1.3 as the primary deployment target.  Hybrid authentication
schemes --- combining a classical signature algorithm such as ECDSA P-256 with a
post-quantum counterpart such as ML-DSA-44 --- provide transitional resilience
against both contemporary and quantum-capable adversaries, but the substantial
difference in signature sizes between classical and post-quantum schemes
introduces fragmentation of TLS handshake messages across multiple TCP segments.
This fragmentation incurs additional network round trips and degrades server
throughput by up to 55\% relative to classical-only deployments.

This paper presents \textit{Deferred Asynchronous PQC Verification} (DAPV), a
protocol-level optimisation that separates the two verification stages of a
hybrid TLS~1.3 handshake.  The classical ECDSA component is verified
synchronously, completing the handshake and deriving session keys immediately;
the post-quantum component is dispatched to a background thread for concurrent
verification.  On PQC success the session is promoted to full authentication;
on failure the connection is immediately terminated.  This architecture removes
the per-connection computational cost of PQC verification from the handshake
critical path while preserving eventual quantum-safe authentication guarantees.

Four contributions are made.  First, a proof-of-concept implementation is
presented within OpenSSL~3.x using the Open Quantum Safe (OQS) provider,
validated on both Windows~11 and Ubuntu~22.04 with AVX2-optimised PQC
primitives.  Second, controlled experiments across RTT values of 0--150\ms\ and
concurrency levels of 100--1000 clients establish that DAPV reduces
client-perceived handshake latency by a constant offset of approximately 20\ms,
independent of RTT, with server throughput improving by 30--36\% for hybrid
P-256\,+\,ML-DSA-44 at 500 concurrent clients.  Third, decomposition of the
Window of Vulnerability (\Dt) --- the interval between provisional handshake
completion and PQC confirmation --- reveals that the measured 14--21\ms\ window
is dominated by thread scheduling overhead rather than cryptographic
computation, which costs less than 0.15\ms\ on modern AVX2 hardware.  Fourth,
a quantitative risk characterisation provides leakage bounds across link speeds
and CPU load levels, identifies a novel denial-of-service vector specific to
deferred verification, and proposes concrete mitigations including a bounded
thread pool, a data-transmission cap, and a minimal OpenSSL API design for
application-layer policy enforcement.
\end{abstract}

\begin{keyword}
Post-Quantum Cryptography \sep TLS 1.3 \sep Hybrid Authentication \sep
Deferred Verification \sep ML-DSA \sep FN-DSA \sep Handshake Latency \sep
Window of Vulnerability
\end{keyword}
\end{frontmatter}

%% ============================================================
\section{Introduction}
\label{sec:intro}

\subsection{Context and Motivation}

Public-key cryptography as deployed in contemporary secure communication
protocols rests on the computational intractability of integer factorisation and
the discrete logarithm problem.  Shor's algorithm~\cite{Shor1994} renders both
problems efficiently solvable on a sufficiently capable quantum computer,
placing the long-term security of protocols such as TLS directly at risk.
While quantum hardware of the scale required to threaten current key sizes does
not yet exist, the threat is operationally significant today because of the
``Harvest Now, Decrypt Later'' (HNDL) attack model~\cite{Alnahawi2024}: an
adversary who records encrypted traffic today may decrypt it retroactively once
quantum capability matures, exposing data with multi-year sensitivity lifetimes.
In recognition of this risk, NIST has standardised a suite of post-quantum
algorithms~\cite{FIPS204}, and national cybersecurity agencies including the
NSA~\cite{NSA2022} and ENISA~\cite{ENISA2023} have issued migration timelines
with near-term compliance obligations.

TLS~1.3~\cite{Cremers2016,Cheval2022} carries the authentication and
confidentiality infrastructure of the modern Internet, protecting web traffic,
API communications, cloud workloads, and mobile ecosystems across billions of
connections per day.  Migrating TLS to post-quantum cryptography is therefore
a matter of immediate operational importance.

\subsection{The Hybrid Transition Strategy}

Direct replacement of classical algorithms with post-quantum alternatives
carries interoperability risk during the transition period.  Hybrid
authentication addresses this by concatenating a classical and a post-quantum
signature within a single TLS~1.3 \texttt{CertificateVerify}
message~\cite{Bindel2017,WiggersIETF2024}, preserving security against
classical adversaries while simultaneously providing quantum-safe authentication
for clients that can verify it.

The performance cost is non-trivial.  The ECDSA P-256 signature occupies
approximately 72~bytes, whereas the NIST-standardised ML-DSA-44
signature~\cite{FIPS204} requires approximately 2420~bytes---a 35-fold
increase.  The combined payload for a P-256\,+\,ML-DSA-44 hybrid handshake
exceeds 3850~bytes, far surpassing the standard Ethernet MTU of 1500~bytes.
TCP fragmentation introduces additional round trips proportional to network
latency, disproportionately penalising wide-area and mobile connections.
Empirical studies consistently report latency increases of 50--300\% and
throughput reductions exceeding 50\%~\cite{Paquin2020,Sikeridis2020a,
Sikeridis2020b,Montenegro2026}.

\subsection{Research Gap}

Existing work on PQC integration into TLS focuses primarily on measurement,
documenting the overhead of candidate algorithms across hardware and network
conditions.  Architectural mitigations remain largely unexplored.  PQC
signature \emph{verification} on modern processors is extremely fast---ML-DSA-44
and FN-DSA-512 complete in well under 1~ms on AVX2-capable
hardware~\cite{OQS2017}---yet verification is scheduled synchronously on the
handshake critical path.  Decoupling PQC verification from the critical path
while preserving its eventual security guarantee is the central design insight
of DAPV.

\subsection{Contributions}

\begin{enumerate}[label=(\arabic*), leftmargin=*]
\item \textbf{Protocol design.}  DAPV is specified as a modification to the
  TLS~1.3 client-side handshake state machine, defining the
  \textsc{provisional} authentication state, conditions for entry and
  promotion, and transcript-binding properties that prevent replay.

\item \textbf{Multi-OS implementation.}  A proof-of-concept is constructed
  within OpenSSL~3.x using the OQS provider, supporting both thread-per-request
  and bounded thread-pool concurrency models, validated on Windows~11 and
  Ubuntu~22.04.  Both testbeds use x86-64 architecture; ARM and embedded
  platforms are a direction for future work.

\item \textbf{Empirical evaluation.}  Controlled experiments across five
  cryptographic configurations, five RTT conditions, and three concurrency
  levels characterise DAPV's latency, throughput, CPU utilisation, and \Dt\
  in detail.  The Window of Vulnerability is decomposed into cryptographic
  and scheduling components, and its behaviour under CPU load is quantified.

\item \textbf{Security characterisation.}  Leakage bounds are derived for the
  provisional window.  A novel denial-of-service attack vector specific to
  deferred verification is identified, and a concrete OpenSSL API for
  application-layer policy enforcement is presented.
\end{enumerate}

\subsection{Scope and Limitations}

DAPV targets the hybrid deployment scenario in which a classical signature
component provides the synchronous provisional anchor.  PQC-only deployments
are outside its current scope.  The provisional window provides classical-level
authentication only; quantum-safe guarantees are established solely upon
completion of background PQC verification.  DAPV is therefore unsuitable for
compliance-critical applications where any interval of unverified session state
is impermissible.

The remainder of this paper is organised as follows.  Section~\ref{sec:related}
surveys background and related work.  Section~\ref{sec:design} presents the
DAPV protocol design, implementation, and analytical framework.
Section~\ref{sec:setup} describes the experimental setup.
Section~\ref{sec:results} reports results.  Section~\ref{sec:security}
provides the security analysis.  Section~\ref{sec:deployment} discusses
deployment considerations.  Section~\ref{sec:comparison} situates results
against prior work.  Section~\ref{sec:conclusion} concludes.

%% ============================================================
\section{Background and Related Work}
\label{sec:related}

\subsection{TLS~1.3 Authentication and Transcript Integrity}

TLS~1.3 provides server authentication through a certificate-based handshake.
The server transmits a \texttt{CertificateVerify} message containing a
signature over the handshake transcript, followed by a \texttt{Finished} MAC.
The \texttt{Finished} MAC is computed over the complete transcript hash, which
includes every preceding handshake message---including the full byte content
of \texttt{CertificateVerify}.  This property is central to DAPV's security:
the PQC signature bytes are present in the transcript and bound into the client
\texttt{Finished} message before any application data is exchanged, even though
their verification is deferred.

\subsection{PQC Performance in TLS}

Paquin, Stebila, and Tamvada~\cite{Paquin2020} conducted the first systematic
benchmarking of post-quantum algorithms in TLS, reporting latency increases of
80--300\% for lattice-based signature schemes in high-RTT environments.
Sikeridis et al.~\cite{Sikeridis2020a,Sikeridis2020b} studied hybrid
authentication, observing throughput reductions exceeding 50\% and attributing
these primarily to TCP fragmentation rather than computational cost.  Montenegro
et al.~\cite{Montenegro2026} confirmed these findings across a broader range of
hybrid configurations.  Open Quantum Safe benchmarking~\cite{OQS2017}
established that ML-DSA-44 and FN-DSA-512 verification on AVX2-capable
hardware completes in 10--100~$\mu$s, confirming that the latency penalty is
dominated by network fragmentation.  Recent work by Sosnowski et
al.~\cite{Sosnowski2023} provides additional characterisation across diverse
network conditions.

\subsection{Hybrid Signature Constructions}

Bindel et al.~\cite{Bindel2017} formalised hybrid certificate and signature
schemes with composable unforgeability guarantees using a nesting construction.
The present implementation uses a simpler concatenation encoding
$\langle\sigma_E\,\|\,\sigma_P\rangle$ consistent with IETF draft
conventions~\cite{WiggersIETF2024}.  This does not carry the same formal
guarantees; production deployments should adopt a construction satisfying the
composable unforgeability definition of Bindel et al.

\subsection{Deferred Authentication}

Post-handshake authentication was proposed for TLS~1.3.  Cremers et
al.~\cite{Cremers2016} identified that without explicit transcript binding, this
mechanism is vulnerable to replay attacks.  Symbolic analyses of TLS~1.3
post-handshake flows~\cite{Cheval2022} confirmed that deferred authentication
requires strict cryptographic context binding.  In DAPV, the PQC signature is
included in the transcript before the client \texttt{Finished} message; the
vulnerability identified by Cremers et al.\ does not apply.

\subsection{Industrial Deployments and Standardisation}

Google's CECPQ1 and CECPQ2 experiments~\cite{Langley2016CECPQ1,
Langley2018CECPQ2,Braithwaite2016} integrated post-quantum key exchange into
Chrome, demonstrating operational feasibility at scale.  Cloudflare deployed
PQC key exchange across CDN edge servers~\cite{Kwiatkowski2019}, reporting
measurable but manageable latency impacts.  The IETF
\texttt{draft-celi-wiggers-tls-authkem}~\cite{WiggersIETF2024} specifies
KEM-based authentication; hybrid signature standardisation is a separate,
ongoing effort.  Neither industrial programme has addressed PQC authentication
latency as a design objective.

%% ============================================================
\section{System Design}
\label{sec:design}

\subsection{Design Objectives}

DAPV is designed to satisfy three objectives simultaneously: (i)~reduce
client-observed handshake latency to near-classical levels; (ii)~preserve the
eventual quantum-safe authentication guarantee of the hybrid scheme; and
(iii)~confine the modification to the client-side verification step, requiring
no changes to the server, wire format, or TLS record layer.

\subsection{Deployment Modes}
\label{sec:modes}

Three deployment configurations are distinguished.  In \textit{classical-only}
operation (ECDSA alone), DAPV is not applicable.  In \textit{hybrid} operation
(classical\,+\,PQC), the classical signature provides a lightweight synchronous
anchor for the provisional state; PQC verification is deferred.  This is the
operative deployment scenario analysed in this paper.  In \textit{PQC-only}
operation (no classical component), DAPV cannot function in its present form
because no lightweight synchronous operation is available to anchor the
provisional state; this is an open design challenge identified in the future
work.

\subsection{Protocol Architecture}
\label{sec:protocol}

Figure~\ref{fig:sequence} illustrates the DAPV-modified TLS~1.3 handshake.
Processing proceeds in six stages.

\begin{enumerate}[leftmargin=*, noitemsep]
\item The client receives \texttt{CertificateVerify} containing the hybrid
  signature $\langle\sigma_E\|\sigma_P\rangle$.

\item The ECDSA component $\sigma_E$ is verified synchronously (tens of
  microseconds), providing the provisional authentication decision.

\item On success, the session transitions to \textsc{provisional}: session
  keys are derived and the client sends its \texttt{Finished} message.
  The \texttt{Finished} MAC covers the complete transcript, including $\sigma_P$
  bytes, establishing cryptographic binding before any application data is sent.

\item $\sigma_P$ is dispatched to a background thread for asynchronous
  verification.

\item On successful PQC verification, the session transitions to
  \textsc{fully\_authenticated}.

\item PQC verification failure at any point terminates the connection with
  \texttt{bad\_certificate}.
\end{enumerate}

\begin{figure}[htbp]
  \centering
  \includegraphics[width=0.92\columnwidth]{dapv_sequence_preview1.pdf}
  \caption{DAPV-modified TLS~1.3 handshake.
    ($\sigma_E$: ECDSA component; $\sigma_P$: PQC component.)
    The client \texttt{Finished} MAC is computed over the full handshake
    transcript including $\sigma_P$ bytes, establishing cryptographic
    binding before any application data is exchanged.  The Window of
    Vulnerability ($\Delta t$) spans the interval between provisional
    completion and PQC confirmation.}
  \label{fig:sequence}
\end{figure}



% \begin{figure}[htbp]
%   \centering
%   \includesvg[width=0.92\columnwidth]{dapv_sequence_preview.svg}
%   \caption{DAPV-modified TLS~1.3 handshake.
%     ($\sigma_E$: ECDSA component; $\sigma_P$: PQC component.)
%     The client \texttt{Finished} MAC is computed over the full handshake
%     transcript including $\sigma_P$ bytes, establishing cryptographic
%     binding before any application data is exchanged.  The Window of
%     Vulnerability ($\Delta t$) spans the interval between provisional
%     completion and PQC confirmation.}
%   \label{fig:sequence}
% \end{figure}

% \begin{figure}[htbp]
%   \centering
%   \includesvg[width=0.92\columnwidth]{dapv_sequence_preview_latex.svg}
%   \caption{DAPV-modified TLS~1.3 handshake.
%     ($\sigma_E$: ECDSA component; $\sigma_P$: PQC component.)
%     The client \texttt{Finished} MAC is computed over the full handshake
%     transcript including $\sigma_P$ bytes, establishing cryptographic
%     binding before any application data is exchanged.  The Window of
%     Vulnerability ($\Delta t$) spans the interval between provisional
%     completion and PQC confirmation.}
%   \label{fig:sequence}
% \end{figure}


\subsection{Analytical Framework}
\label{sec:analytical}

\subsubsection{Handshake Latency Model}

Classical TLS~1.3 handshake latency follows an affine function of RTT:
\begin{equation}
  L_\text{classical} = L_0 + k \cdot \text{RTT}
  \label{eq:classical}
\end{equation}
where $L_0$ is the base computational latency and $k$ is the number of
required handshake round trips.  Under hybrid PQC authentication, enlarged
payloads trigger TCP fragmentation, introducing $\delta k \geq 0$ additional
round trips:
\begin{equation}
  L_\text{hybrid} = L_0 + (k + \delta k) \cdot \text{RTT}
  \label{eq:hybrid}
\end{equation}
DAPV removes PQC verification time $T_\text{verify}$ from the critical path
but does not alter fragmentation behaviour:
\begin{equation}
  L_\text{DAPV} \approx L_0 + (k + \delta k) \cdot \text{RTT} - T_\text{verify}
  \label{eq:dapv}
\end{equation}
Because $T_\text{verify}$ is RTT-independent, the latency benefit of DAPV is a
\emph{constant offset} relative to the synchronous hybrid curve, not a change
in slope.

\subsubsection{Throughput Model}

Server throughput under hybrid authentication is modulated by PQC signature
size $S_\text{pqc}$ relative to the MTU $S_\text{mtu} = 1500$~bytes:
\begin{equation}
  T_\text{hybrid} = \frac{T_\text{classical}}
    {1 + \lambda \cdot S_\text{pqc} / S_\text{mtu}}
\end{equation}
where $\lambda \geq 0$ captures RTT and congestion effects.

\subsubsection{Leakage and Risk}

The maximum application data traversing the provisional channel during \Dt\ is:
\begin{equation}
  L_\text{max} = \frac{B \times \Dt}{8}
  \label{eq:leakage}
\end{equation}
where $B$ is link bandwidth in bits per second.  A data-transmission cap
$D_\text{policy}$ applied at the application layer bounds this irrespective of
\Dt\ and $B$.  The provisional exposure probability follows:
\begin{equation}
  R(\Dt, B) = 1 - e^{-\beta B \Dt}
  \label{eq:risk}
\end{equation}
where $\beta$ is an empirical scaling factor calibrated to the acceptable
risk threshold.

\subsubsection{System Efficiency}

\begin{equation}
  E_\text{DAPV} = \frac{T_\text{DAPV} / T_\text{classical}}
                       {L_\text{DAPV} / L_\text{hybrid}}
  \label{eq:efficiency}
\end{equation}
$E_\text{DAPV} > 1$ indicates a net system-level benefit: proportionally greater
throughput improvement than residual latency cost.

\subsection{Algorithms}
\label{sec:algorithms}

Algorithms~\ref{alg:client}--\ref{alg:policy} formalise the client handshake
procedure, server dispatch mechanism, and application-layer policy enforcement.

\begin{algorithm}[H]
\caption{DAPV Client Handshake Procedure}
\label{alg:client}
\DontPrintSemicolon
\KwIn{Hybrid signature $\langle\sigma_E, \sigma_P\rangle$ in \texttt{CertificateVerify}}
\KwOut{Session state $\in$ \{\textsc{provisional}, \textsc{fully\_authenticated}\}}
\BlankLine
Extract $\sigma_E$ and $\sigma_P$ from \texttt{CertificateVerify}\;
\uIf{$\mathrm{VerifyECDSA}(\sigma_E,\;\mathrm{cert}) = \mathrm{true}$}{
  $\mathrm{session.state} \leftarrow \textsc{provisional}$\;
  Derive session keys\;
  Compute \texttt{Finished} MAC over full transcript (includes $\sigma_P$ bytes)\;
  Send \texttt{Finished}\;
  Dispatch $\mathrm{AsyncVerify}(\sigma_P,\;\mathrm{cert},\;\mathrm{sid})$
    to background thread\;
}
\Else{
  Send \texttt{bad\_certificate}\;  Abort connection\;
}
\BlankLine
\SetKwProg{Proc}{procedure}{:}{}
\Proc{$\mathrm{AsyncVerify}(\sigma_P,\;\mathrm{cert},\;\mathrm{sid})$}{
  \uIf{$\mathrm{VerifyPQC}(\sigma_P,\;\mathrm{cert}) = \mathrm{success}$}{
    $\mathrm{session}(\mathrm{sid}).\mathrm{state}
      \leftarrow \textsc{fully\_authenticated}$\;
  }
  \Else{
    Terminate $\mathrm{session}(\mathrm{sid})$\;
    Send \texttt{bad\_certificate}\;
  }
}
\end{algorithm}

\begin{algorithm}[H]
\caption{Bounded Server Thread-Pool Dispatcher}
\label{alg:server}
\DontPrintSemicolon
\KwIn{Verification task queue; pool bound $M$}
Initialise thread pool $\mathcal{P}$ with $M$ workers\tcp*{$M$ bounded
  to resist provisional flooding}
\While{server running}{
  \If{$\mathrm{TaskQueue} \neq \emptyset$}{
    $\mathrm{task} \leftarrow \mathrm{Dequeue}(\mathrm{TaskQueue})$\;
    Assign task to idle worker $w \in \mathcal{P}$\;
    $w.\mathrm{execute}\bigl(
      \mathrm{VerifyPQC}(\mathrm{task}.\sigma_P,\;\mathrm{task.cert})
    \bigr)$\;
  }
}
\end{algorithm}

\begin{algorithm}[H]
\caption{Application-Layer Data Transmission Policy}
\label{alg:policy}
\DontPrintSemicolon
\KwIn{\texttt{session.state}; outbound payload size $d$; cap $D_\text{policy}$}
\uIf{$\mathrm{session.state} = \textsc{provisional}$}{
  \uIf{$d > D_\text{policy}$}{
    Buffer payload\;
    Retry after \textsc{fully\_authenticated}\;
  }
  \Else{
    Transmit payload\;
  }
}
\Else{
  Transmit payload unconditionally\;
}
\end{algorithm}

\subsection{Implementation}
\label{sec:impl}

The proof-of-concept is built on OpenSSL~3.2 with the OQS provider and
liboqs~\cite{OQS2017}, patching the \texttt{oqs\_sig\_verify} callback to
dispatch the PQC component asynchronously via \texttt{CreateThread} (Windows)
or \texttt{pthread} (Linux).  Hybrid certificates concatenate ECDSA P-256 and
PQC public keys; signatures are encoded as
$\langle\mathrm{len}(\sigma_E)\|\sigma_E\|\sigma_P\rangle$.

Three concurrency models are supported: (i)~\textit{thread per request}, used
in this prototype for simplicity; (ii)~\textit{bounded thread pool}, recommended
for production to cap resource consumption and resist denial-of-service;
and (iii)~\textit{event-driven dispatch}, suitable for high-concurrency edge
deployments.

%% ============================================================
\section{Experimental Setup}
\label{sec:setup}

\subsection{Hardware Platforms}

Experiments were conducted on two testbeds.  The Windows testbed comprised an
Intel Core i7-11800H (8 cores, 2.3~GHz, 32~GB RAM) running Windows~11 Pro.
The Linux testbed comprised an Intel Xeon Silver~4310 (12 cores, 2.1~GHz,
64~GB RAM) running Ubuntu~22.04 LTS.  Both systems used dedicated Gigabit
Ethernet.  AVX2 support was confirmed active on both platforms prior to
benchmarking.  Both testbeds share x86-64 microarchitecture; evaluation on
ARM-class embedded hardware is a direction for future work.

\subsection{Software Stack}

OpenSSL~3.2 was compiled with the OQS provider~\cite{OQS2017} and liboqs,
providing AVX2-optimised implementations of all PQC schemes under test.
Compiler toolchains were Visual Studio~2022 (Windows) and GCC~11.3 (Linux).
DAPV modifications were confined to the \texttt{oqs\_sig\_verify} callback and
connection state machine; no changes were required to TLS record processing or
key derivation.

\subsection{Cryptographic Configurations}

Five configurations were benchmarked:
\begin{itemize}[noitemsep]
  \item \textbf{ECDSA P-256} --- classical baseline.
  \item \textbf{FN-DSA-512} --- NIST Level~1; FIPS~206 draft at time of
        experiments; compact signatures with fast verification.
  \item \textbf{ML-DSA-44} --- NIST Level~2; FIPS~204 final standard; larger
        signatures with well-studied lattice security.
  \item \textbf{P-256\,+\,FN-DSA-512} --- hybrid; primary DAPV target.
  \item \textbf{P-256\,+\,ML-DSA-44} --- hybrid; primary DAPV target.
\end{itemize}

\noindent Standalone FN-DSA-512 and ML-DSA-44 DAPV measurements (without a
classical anchor) are included to characterise threading overhead; they do not
represent operative DAPV deployments and should be interpreted accordingly.

\subsection{Network Emulation}

One-way delay was introduced using Clumsy~0.3 (Windows) and \texttt{tc netem}
(Linux), producing RTT values of 0, 20, 50, 100, and 150\ms.  Packet loss was
held at zero to isolate fragmentation effects from retransmission events.

\subsection{Measurement Procedure}

Four metrics were collected across 30 repetitions per configuration:

\begin{enumerate}[noitemsep]
\item \textbf{Handshake latency} (ms): wall-clock interval from
  \texttt{ClientHello} transmission to receipt of server \texttt{Finished},
  measured at the client.

\item \textbf{Server throughput} (conn/s): maximum sustained TLS handshake
  rate under 100, 500, and 1000 concurrent clients.

\item \textbf{Peak CPU utilisation} (\%): server-side peak during handshake
  bursts.

\item \textbf{Window of Vulnerability (\Dt)}: decomposed into (a)~raw PQC
  cryptographic verification time (timed directly without thread dispatch) and
  (b)~total wall-clock \Dt\ including thread creation and scheduling.
\end{enumerate}

\noindent Standard deviation was below 4\% of the mean in all cases;
95\% confidence intervals are reported throughout.

%% ============================================================
\section{Results}
\label{sec:results}

\subsection{Handshake Latency at Zero RTT}
\label{sec:res_latency0}

Figure~\ref{fig:latency_bar} presents mean handshake latency at 0\ms\ RTT,
where differences reflect computational overhead alone.  DAPV reduces latency
for both hybrid configurations: from 48.5 to 34.2\ms\ for
P-256\,+\,FN-DSA-512 (a reduction of 29\%) and from 55.6 to 34.9\ms\ for
P-256\,+\,ML-DSA-44 (37\%).  The classical P-256 baseline is unaffected.

\begin{figure}[htbp]
  \centering
  \includegraphics[width=0.92\columnwidth]{fig01_latency_bar_0ms.pdf}
  \caption{Mean handshake latency at 0\ms\ RTT for all five configurations
    under synchronous (blue) and DAPV (red) scheduling.  Green arrows indicate
    the absolute and percentage reduction for the two operative hybrid
    configurations (shaded background).  Standalone PQC DAPV entries
    quantify threading overhead only.}
  \label{fig:latency_bar}
\end{figure}

\subsection{Latency as a Function of RTT}
\label{sec:res_latency_rtt}

Figure~\ref{fig:latency_rtt} presents handshake latency versus RTT for
P-256\,+\,ML-DSA-44.  All three curves grow linearly with RTT, consistent with
Equations~\eqref{eq:classical} and~\eqref{eq:hybrid}: the dominant driver at
high RTT is the fragmentation-induced extra round trip.  DAPV maintains a
constant offset of approximately 20\ms\ below the synchronous hybrid curve
across the full tested range, confirming the RTT-independence of
$T_\text{verify}$ in Equation~\eqref{eq:dapv}.

\begin{figure}[htbp]
  \centering
  \includegraphics[width=0.92\columnwidth]{fig02_latency_vs_rtt.pdf}
  \caption{Handshake latency versus RTT for P-256\,+\,ML-DSA-44.  DAPV
    delivers a constant $\approx$20\ms\ reduction relative to synchronous
    hybrid at every tested RTT value.  The double-headed arrow at RTT\,=\,150\ms\
    marks this constant offset.  Both DAPV and synchronous hybrid curves have
    equal slopes, confirming that fragmentation-induced extra round trips are
    not eliminated by DAPV.}
  \label{fig:latency_rtt}
\end{figure}

Figure~\ref{fig:latency_saving} isolates the absolute latency saving provided
by DAPV for each configuration at three RTT values.  Savings vary by less than
0.2\ms\ across RTT conditions, confirming the RTT-independence of
$T_\text{verify}$.  P-256\,+\,ML-DSA-44 yields the largest saving (20.7\ms)
because ML-DSA-44 carries a higher per-call computational cost than FN-DSA-512.

\begin{figure}[htbp]
  \centering
  \includegraphics[width=0.92\columnwidth]{fig03_latency_saving_by_rtt.pdf}
  \caption{Absolute latency reduction (Sync\,$-$\,DAPV) for each configuration
    at RTT\,=\,0, 50, and 150\ms.  The near-constant saving across RTT
    conditions---variation below 0.2\ms---confirms that $T_\text{verify}$ is
    RTT-independent.  P-256\,+\,ML-DSA-44 achieves the largest reduction
    owing to ML-DSA-44's higher computational cost.}
  \label{fig:latency_saving}
\end{figure}

\subsection{Server Throughput}
\label{sec:res_throughput}

Figure~\ref{fig:throughput_bar} presents server throughput at 100 and 1000
concurrent clients.  Classical ECDSA sustains approximately 820~conn/s.  Hybrid
schemes degrade to 390--420~conn/s---a reduction exceeding 50\%.  DAPV
restores throughput to 530--565~conn/s at 100 clients (30--36\% improvement
over synchronous hybrid).  The improvement arises from earlier release of
server-side connection state: once the client sends its \texttt{Finished}
message under provisional authentication, the server may begin processing the
next connection without waiting for PQC verification.

\begin{figure}[htbp]
  \centering
  \includegraphics[width=0.92\columnwidth]{fig04_throughput_bar.pdf}
  \caption{Server throughput at 100 (solid bars) and 1000 (hatched bars)
    concurrent clients.  Blue bars denote synchronous operation; red bars
    denote DAPV.  DAPV improves throughput by 30--36\% for hybrid
    P-256\,+\,ML-DSA-44 at 100 clients, declining to $\approx$30\% at 1000
    as thread-pool scheduling contention increases.
    Standard deviation below 4\% of mean throughout.}
  \label{fig:throughput_bar}
\end{figure}

Figure~\ref{fig:throughput_rtt} extends this analysis across RTT values at
500 concurrent clients.  Both synchronous hybrid and DAPV hybrid degrade
monotonically with RTT as fragmentation extends connection dwell time.  DAPV
maintains a 30--34\% throughput advantage across the full RTT range.

\begin{figure}[htbp]
  \centering
  \includegraphics[width=0.92\columnwidth]{fig05_throughput_vs_rtt.pdf}
  \caption{Server throughput versus RTT at 500 concurrent clients.
    The shaded region between synchronous and DAPV hybrid curves represents
    the DAPV throughput gain.  Percentage labels indicate the gain at each
    RTT value.  Both curves degrade with RTT due to fragmentation; the
    relative gain is stable at 30--34\%.}
  \label{fig:throughput_rtt}
\end{figure}

Figure~\ref{fig:uplift} summarises throughput uplift across all tested
concurrency levels.  Uplift is stable at 34--36\% for P-256\,+\,ML-DSA-44 and
31--35\% for P-256\,+\,FN-DSA-512, confirming that the improvement is robust
to concurrency within the tested range.

\begin{figure}[htbp]
  \centering
  \includegraphics[width=0.82\columnwidth]{fig06_throughput_uplift.pdf}
  \caption{DAPV throughput uplift (\%) relative to synchronous hybrid, at
    100, 500, and 1000 concurrent clients.  The shaded band illustrates the
    stable improvement region.  Uplift is consistent across concurrency levels,
    suggesting the benefit is not sensitive to client load within this range.}
  \label{fig:uplift}
\end{figure}

\subsection{CPU Utilisation}
\label{sec:res_cpu}

Figure~\ref{fig:cpu} presents peak server CPU utilisation during handshake
bursts.  Classical ECDSA peaks at 35\%; hybrid configurations peak at
79--85\%.  DAPV does not reduce aggregate CPU consumption; the same
cryptographic operations execute in both scheduling models.  What DAPV changes
is the \textit{temporal distribution} of that work: rather than sustained
high-utilisation plateaus coinciding with client-observed handshake completion,
computation is distributed into background thread activity that overlaps with
the handling of subsequent connections.  This redistribution reduces peak
utilisation episodes and improves scheduler responsiveness under high
concurrency, contributing to the throughput improvement reported above.

\begin{figure}[htbp]
  \centering
  \includegraphics[width=0.88\columnwidth]{fig07_cpu_utilisation.pdf}
  \caption{Peak server CPU utilisation during handshake bursts (synchronous
    configurations).  Bars are colour-coded by utilisation tier: green\,=\,low
    ($<$50\%), orange\,=\,moderate (50--80\%), red\,=\,high ($>$80\%).
    DAPV redistributes PQC computation into background threads, reducing peak
    episodes without changing aggregate work.
    95\% confidence intervals within $\pm$3 percentage points.}
  \label{fig:cpu}
\end{figure}

\subsection{Window of Vulnerability Decomposition}
\label{sec:res_delta_t}

Figure~\ref{fig:deltat_decomp} decomposes the total \Dt\ into its constituent
components.  The PQC cryptographic verification operation completes in
0.08--0.14\ms\ on AVX2 hardware, consistent with published OQS benchmarks.
Thread creation, OS scheduling, and context-switching overhead account for the
remainder, contributing 14--21\ms\ to the total prototype \Dt.

\begin{figure}[htbp]
  \centering
  \includegraphics[width=0.92\columnwidth]{fig08_delta_t_decomposition.pdf}
  \caption{\Dt\ decomposed into PQC cryptographic verification time (cyan,
    $<$0.15\ms) and thread dispatch and OS scheduling overhead (purple,
    14--21\ms).  The cryptographic component accounts for less than 1\% of
    the total \Dt\ in the prototype implementation.  An optimised bounded
    thread pool would bring total \Dt\ close to the cryptographic bound.}
  \label{fig:deltat_decomp}
\end{figure}

This decomposition has an important security implication: the true exposure
window --- the interval during which a forged message could go undetected ---
corresponds to the cryptographic verification time (below 0.15\ms) rather than
to the implementation-overhead \Dt\ (14--21\ms).

Table~\ref{tab:delta_t} provides full statistical characterisation of \Dt\
across 30 trials per configuration.

\begin{table}[htbp]
\centering
\caption{Window of Vulnerability (\Dt) decomposition, 30 trials per
  configuration, idle hardware.  All values in ms.}
\label{tab:delta_t}
\small
\begin{tabular}{lrrrr}
\toprule
\textbf{Scheme} & \textbf{Crypto} & \textbf{Thread} &
\textbf{Total \Dt} & \textbf{95\% CI} \\
\midrule
FN-DSA-512            & 0.12 & 18.28 & 18.4 & [16.6,\;20.2] \\
ML-DSA-44             & 0.08 & 14.42 & 14.5 & [13.1,\;15.9] \\
P-256\,+\,FN-DSA-512  & 0.14 & 18.76 & 18.9 & [16.7,\;21.1] \\
P-256\,+\,ML-DSA-44   & 0.10 & 21.10 & 21.2 & [18.6,\;23.8] \\
\bottomrule
\end{tabular}
\end{table}

Figure~\ref{fig:deltat_load} characterises \Dt\ under server CPU load for
P-256\,+\,ML-DSA-44.  At 80\% CPU utilisation, the 95th-percentile \Dt\
reaches 58.7\ms---approximately 2.5 times the idle-hardware median---driven by
OS scheduling latency for background threads competing with foreground
connection processing.  Deployment policy must therefore calibrate the
data-transmission cap $D_\text{policy}$ against worst-case \Dt\ under
production load rather than idle-hardware measurements.  Based on these
results, a conservative cap of $D_\text{policy} \leq 8$~KB is recommended.

\begin{figure}[htbp]
  \centering
  \includegraphics[width=0.88\columnwidth]{fig09_delta_t_vs_cpu_load.pdf}
  \caption{\Dt\ for P-256\,+\,ML-DSA-44 at three server CPU load levels.
    Each load level shows median, mean, and 95th-percentile \Dt\ across
    30 trials.  At 80\% CPU load, the 95th-percentile reaches 58.7\ms,
    approximately 2.5$\times$ the idle median, due to OS scheduling contention.
    The dashed line marks the idle-hardware 95th percentile for reference.}
  \label{fig:deltat_load}
\end{figure}

\subsection{Signature Size and Fragmentation}
\label{sec:res_size}

Figure~\ref{fig:overhead} presents the combined public key and signature payload
alongside the 1500-byte Ethernet MTU boundary.  FN-DSA-512 at 1587~bytes
marginally exceeds the MTU, requiring two TCP segments.  ML-DSA-44 at
2420~bytes requires three.  The hybrid P-256\,+\,ML-DSA-44 combination at
3858~bytes requires at minimum three segments, commonly four or more with
certificate overhead.  These figures explain the $\delta k$ fragmentation
penalty in Equation~\eqref{eq:hybrid} and confirm that DAPV's 20\ms\ latency
improvement cannot eliminate the extra round trips induced by fragmentation.

\begin{figure}[htbp]
  \centering
  \includegraphics[width=0.88\columnwidth]{fig10_signature_size.pdf}
  \caption{Combined public key and signature payload per configuration.
    The dashed red line marks the 1500-byte Ethernet MTU.  Green bars fall
    within the MTU; red bars require TCP fragmentation.  The segment count
    for each configuration is annotated.  DAPV defers verification but
    cannot reduce payload size; fragmentation persists in all hybrid
    configurations.}
  \label{fig:overhead}
\end{figure}

\subsection{Consolidated Results}
\label{sec:res_summary}

Table~\ref{tab:summary} presents a unified comparison across all evaluated
configurations at 500 concurrent clients.

\begin{table}[htbp]
\centering
\caption{Consolidated performance comparison across all configurations.
  Throughput at 500 concurrent clients; \Dt\ on Windows testbed.
  \textbf{Bold}: operative DAPV results.
  $\dagger$: standalone PQC with threaded dispatch; no classical provisional
  anchor; included as threading-overhead reference only
  (see Section~\ref{sec:modes}).}
\label{tab:summary}
\small
\setlength{\tabcolsep}{6pt}
\resizebox{\columnwidth}{!}{%
\begin{tabular}{lrrrrr}
\toprule
\textbf{Configuration} &
  \textbf{Latency 0\ms} & \textbf{Latency 150\ms} &
  \textbf{Throughput} & \textbf{Peak CPU} & \textbf{\Dt} \\
& (ms) & (ms) & (conn/s) & (\%) & (ms) \\
\midrule
ECDSA P-256 (classical)               & 15.2 & 165 & 770 & 35 & --- \\[3pt]
FN-DSA-512, synchronous               & 30.1 & 180 & 380 & 76 & --- \\
FN-DSA-512, DAPV$^\dagger$            & 26.7 & 163 & 495 & 76 & 18.4 \\[3pt]
ML-DSA-44, synchronous                & 41.2 & 191 & 390 & 82 & --- \\
ML-DSA-44, DAPV$^\dagger$             & 29.3 & 172 & 510 & 82 & 14.5 \\[3pt]
P-256\,+\,FN-DSA-512, synchronous     & 48.5 & 199 & 440 & 85 & --- \\
\textbf{P-256\,+\,FN-DSA-512, DAPV}   & \textbf{34.2} & \textbf{179} & \textbf{575 (+31\%)} & \textbf{85} & \textbf{18.9} \\[3pt]
P-256\,+\,ML-DSA-44, synchronous      & 55.6 & 206 & 470 & 79 & --- \\
\textbf{P-256\,+\,ML-DSA-44, DAPV}    & \textbf{34.9} & \textbf{185} & \textbf{615 (+31\%)} & \textbf{79} & \textbf{21.2} \\
\bottomrule
\end{tabular}%
}
\end{table}

Table~\ref{tab:efficiency} reports the system efficiency metric
$E_\text{DAPV}$ (Equation~\eqref{eq:efficiency}) for the two hybrid
configurations.  Both exceed 1.0, confirming a net system-level benefit across
both throughput and latency dimensions.

\begin{table}[htbp]
\centering
\caption{System efficiency metric $E_\text{DAPV}$ for hybrid configurations.
  Computed at 500 concurrent clients, 0\ms\ RTT.
  $E_\text{DAPV} > 1$: net system-level improvement.}
\label{tab:efficiency}
\small
\begin{tabular}{lrrrr}
\toprule
\textbf{Scheme} &
  $T_\text{DAPV}/T_\text{cl.}$ &
  $L_\text{DAPV}/L_\text{hyb.}$ &
  $E_\text{DAPV}$ & \textbf{Outcome} \\
\midrule
P-256\,+\,FN-DSA-512 & 0.747 & 0.706 & 1.058 & Net gain \\
P-256\,+\,ML-DSA-44  & 0.799 & 0.628 & 1.272 & Net gain \\
\bottomrule
\end{tabular}
\end{table}

Table~\ref{tab:full_rtt} extends the consolidated view to all five RTT
conditions, presenting complete latency and throughput data for every
configuration.  Latency at intermediate RTT values is obtained by linear
interpolation between the measured 0\ms\ and 150\ms\ endpoints, consistent
with the affine model of Equation~\eqref{eq:latency_dapv}.  Throughput values
for P-256\,+\,ML-DSA-44 across RTT are directly measured
(Figure~\ref{fig:tput_rtt}); those for P-256\,+\,FN-DSA-512 at non-zero RTT
are scaled by the same proportional degradation factor observed for
P-256\,+\,ML-DSA-44 and are marked~$^*$.  Throughput for standalone PQC
configurations was not measured beyond RTT\,=\,0\ms\ and is omitted.

\begin{table*}[htbp]
\centering
\caption{Complete latency and throughput results for all configurations and
  all five tested RTT conditions (0--150\,ms).
  Latency at RTT\,=\,20, 50, 100\,ms is linearly interpolated between the
  measured 0\,ms and 150\,ms endpoints (slope $\approx$\,1\,ms/ms\,RTT for
  all schemes; see Section~\ref{sec:res_latency_rtt}).
  Throughput at 500 concurrent clients; boldface rows are operative DAPV
  deployments.
  ${}^\dagger$~standalone PQC; no classical anchor; throughput measured at
  RTT\,=\,0\,ms only.
  ${}^*$~extrapolated via the RTT-degradation ratio of P-256\,+\,ML-DSA-44.}
\label{tab:full_rtt}
\small
\setlength{\tabcolsep}{4pt}
\resizebox{\textwidth}{!}{%
\begin{tabular}{@{}l r rr rr  rr rr@{}}
\toprule
\multirow{2}{*}{\textbf{Configuration}} &
\multirow{2}{*}{\textbf{RTT}} &
\multicolumn{4}{c}{\textbf{Handshake Latency (ms)}} &
\multicolumn{4}{c}{\textbf{Throughput @ 500 clients (conn/s)}} \\
\cmidrule(lr){3-6}\cmidrule(lr){7-10}
& \textbf{(ms)} &
  \textit{Sync} & \textit{DAPV} & $\Delta L$ (ms) & $\Delta L$ (\%) &
  \textit{Sync} & \textit{DAPV} & Gain (c/s) & Gain (\%) \\
\midrule
ECDSA P-256           &   0 &  15.2 & ---   & ---  & ---  & 770 & --- & ---      & ---    \\
(classical baseline)  &  20 &  35.2 & ---   & ---  & ---  & 710 & --- & ---      & ---    \\
                      &  50 &  65.1 & ---   & ---  & ---  & 640 & --- & ---      & ---    \\
                      & 100 & 115.1 & ---   & ---  & ---  & 570 & --- & ---      & ---    \\
                      & 150 & 165.0 & ---   & ---  & ---  & 495 & --- & ---      & ---    \\
\addlinespace[3pt]
FN-DSA-512            &   0 &  30.1 &  26.7 &  3.4 & 11.3 & 380 & 495 & $+$115  & $+$30.3 \\
(DAPV$^\dagger$)      &  20 &  50.1 &  44.9 &  5.2 & 10.4 & --- & --- & ---     & ---    \\
                      &  50 &  80.1 &  72.2 &  7.9 &  9.9 & --- & --- & ---     & ---    \\
                      & 100 & 130.1 & 117.6 & 12.5 &  9.6 & --- & --- & ---     & ---    \\
                      & 150 & 180.0 & 163.0 & 17.0 &  9.4 & --- & --- & ---     & ---    \\
\addlinespace[3pt]
ML-DSA-44             &   0 &  41.2 &  29.3 & 11.9 & 28.9 & 390 & 510 & $+$120  & $+$30.8 \\
(DAPV$^\dagger$)      &  20 &  61.2 &  48.3 & 12.9 & 21.1 & --- & --- & ---     & ---    \\
                      &  50 &  91.1 &  76.9 & 14.2 & 15.6 & --- & --- & ---     & ---    \\
                      & 100 & 141.1 & 124.4 & 16.7 & 11.8 & --- & --- & ---     & ---    \\
                      & 150 & 191.0 & 172.0 & 19.0 &  9.9 & --- & --- & ---     & ---    \\
\addlinespace[5pt]
\textbf{P-256\,+\,FN-DSA-512}  &   \textbf{0} &  \textbf{48.5} & \textbf{34.2} & \textbf{14.3} & \textbf{29.5} & \textbf{440} & \textbf{575}       & $\mathbf{+135}$ & $\mathbf{+30.7}$ \\
\textbf{(operative DAPV)}       &  \textbf{20} &  \textbf{68.6} & \textbf{53.5} & \textbf{15.1} & \textbf{22.0} & \textbf{412}$^*$ & \textbf{538}$^*$ & $\mathbf{+126}$ & $\mathbf{+30.6}$ \\
                                &  \textbf{50} &  \textbf{98.7} & \textbf{82.5} & \textbf{16.2} & \textbf{16.4} & \textbf{370}$^*$ & \textbf{483}$^*$ & $\mathbf{+113}$ & $\mathbf{+30.5}$ \\
                                & \textbf{100} & \textbf{148.8} & \textbf{130.7} & \textbf{18.1} & \textbf{12.2} & \textbf{332}$^*$ & \textbf{434}$^*$ & $\mathbf{+102}$ & $\mathbf{+30.7}$ \\
                                & \textbf{150} & \textbf{199.0} & \textbf{179.0} & \textbf{20.0} & \textbf{10.1} & \textbf{290}$^*$ & \textbf{379}$^*$ & $\mathbf{+89}$  & $\mathbf{+30.7}$ \\
\addlinespace[5pt]
\textbf{P-256\,+\,ML-DSA-44}   &   \textbf{0} &  \textbf{55.6} & \textbf{34.9} & \textbf{20.7} & \textbf{37.2} & \textbf{470} & \textbf{615} & $\mathbf{+145}$ & $\mathbf{+30.9}$ \\
\textbf{(operative DAPV)}       &  \textbf{20} &  \textbf{75.7} & \textbf{54.9} & \textbf{20.8} & \textbf{27.5} & \textbf{440} & \textbf{580} & $\mathbf{+140}$ & $\mathbf{+31.8}$ \\
                                &  \textbf{50} & \textbf{105.8} & \textbf{85.0} & \textbf{20.8} & \textbf{19.7} & \textbf{395} & \textbf{525} & $\mathbf{+130}$ & $\mathbf{+32.9}$ \\
                                & \textbf{100} & \textbf{155.9} & \textbf{135.0} & \textbf{20.9} & \textbf{13.4} & \textbf{355} & \textbf{470} & $\mathbf{+115}$ & $\mathbf{+32.4}$ \\
                                & \textbf{150} & \textbf{206.0} & \textbf{185.0} & \textbf{21.0} & \textbf{10.2} & \textbf{310} & \textbf{415} & $\mathbf{+105}$ & $\mathbf{+33.9}$ \\
\bottomrule
\end{tabular}%
}
\end{table*}

Three observations emerge from Table~\ref{tab:full_rtt}.
First, the absolute latency saving $\Delta L$ grows with RTT for all
configurations: for P-256\,+\,ML-DSA-44 it rises from 20.7\ms\ at RTT\,=\,0
to 21.0\ms\ at RTT\,=\,150\ms, consistent with Figure~\ref{fig:latency_saving}.
Second, the \emph{relative} saving $\Delta L$ (\%) declines sharply with RTT
--- from 37.2\% to 10.2\% for P-256\,+\,ML-DSA-44 --- because the absolute
saving is fixed while total latency grows linearly; this is an expected
consequence of the RTT-affine model.
Third, the throughput gain for P-256\,+\,ML-DSA-44 is stable within
3 percentage points across the full tested RTT range (30.9--33.9\%),
confirming that DAPV's benefit does not depend on any particular network
condition.

%% ============================================================
\section{Security Analysis}
\label{sec:security}

The analysis presented here provides quantitative risk characterisation,
including leakage bounds, attack taxonomy, and mitigation strategies.  A formal
game-based cryptographic security proof --- for example, an EU-CMA analysis of
the DAPV-modified handshake under combined classical and quantum adversarial
models using Tamarin or ProVerif --- is outside the scope of the present work
and is identified as a priority direction for future research.

\subsection{Threat Model}

The analysis follows the extended Dolev--Yao model: the adversary controls the
communication channel and may intercept, delay, replay, and modify messages.
Two adversary classes are distinguished.

A \textit{classical adversary} cannot break ECDSA in real time.  This is the
relevant adversary during the provisional window \Dt: the ECDSA signature serves
as the synchronous authentication anchor, and a classical adversary cannot
forge $\sigma_E$ to impersonate the server.

A \textit{CRQC-equipped adversary} can in principle execute Shor's algorithm to
forge ECDSA signatures.  This adversary class is outside the threat model of
the provisional window: full quantum-safe authentication is established only
after background PQC verification succeeds.  This scoping is operationally
appropriate under current conditions, as quantum hardware capable of attacking
ECDSA at Internet scale does not currently exist.  As the threat landscape
evolves, the acceptable magnitude of \Dt\ must be re-evaluated.

\subsection{Leakage Bounds}

Table~\ref{tab:leakage} presents provisional data exposure bounds derived from
Equation~\eqref{eq:leakage} for representative link speeds.  Under the
prototype \Dt\ of 14--21\ms, gigabit links can transmit 1.8--2.5~MB during the
provisional window, motivating a conservative $D_\text{policy}$ cap.  With an
optimised thread-pool implementation reducing total \Dt\ below 1\ms, the same
link transmits at most 125~KB, comfortably within a 256-KB cap.

\begin{table}[htbp]
\centering
\caption{Maximum provisionally transmitted data during \Dt.
  Bounds from $L_\text{max} = B \times \Dt / 8$.
  The optimised column assumes a thread-pool implementation reducing total
  \Dt\ below 1\ms.}
\label{tab:leakage}
\small
\begin{tabular}{lrrrr}
\toprule
& \multicolumn{2}{c}{\textbf{Prototype \Dt}} &
  \multicolumn{2}{c}{\textbf{Optimised \Dt}} \\
\cmidrule(lr){2-3}\cmidrule(lr){4-5}
\textbf{Link speed} & @15\ms & @21\ms & @0.5\ms & @1\ms \\
\midrule
1\,Mbps   & 1.8\,KB  & 2.6\,KB  & 62\,B    & 125\,B  \\
100\,Mbps & 183\,KB  & 256\,KB  & 6.1\,KB  & 12.5\,KB \\
1\,Gbps   & 1.8\,MB  & 2.5\,MB  & 61\,KB   & 125\,KB  \\
10\,Gbps  & 17.9\,MB & 25\,MB   & 610\,KB  & 1.25\,MB \\
\bottomrule
\end{tabular}
\end{table}

Figure~\ref{fig:risk} presents the provisional exposure probability
$R(\Dt, B)$ (Equation~\eqref{eq:risk}) as a function of \Dt\ for four link
speeds.  At prototype \Dt\ on gigabit links, $R$ approaches 10\%; with
optimised \Dt\ below 1\ms, $R$ falls below 0.1\% at all tested speeds.

\begin{figure}[htbp]
  \centering
  \includegraphics[width=0.92\columnwidth]{fig11_risk_vs_delta_t.pdf}
  \caption{Provisional exposure risk $R(\Dt, B)$ versus \Dt\ for four link
    speeds.  Vertical dashed lines mark the prototype \Dt\ ($\approx$20\ms)
    and the projected optimised \Dt\ ($<$1\ms).  At prototype \Dt\ on gigabit
    links, risk approaches 10\%; an optimised implementation reduces this
    below 0.1\% at all tested speeds.}
  \label{fig:risk}
\end{figure}

\subsection{Device Scaling}

Table~\ref{tab:scaling} projects \Dt\ for device classes beyond the x86-64
testbeds, based on published OQS benchmarks~\cite{OQS2017}.  DAPV is not
recommended for Cortex-M4 class microcontrollers without dedicated PQC hardware
acceleration, as the resulting \Dt\ of up to 700\ms\ would permit data volumes
far exceeding any practical $D_\text{policy}$ cap.

\begin{table}[htbp]
\centering
\caption{Projected \Dt\ scaling across device classes.
  Cryptographic verification times from OQS benchmarks~\cite{OQS2017};
  thread scheduling estimates extrapolated from prototype measurements.}
\label{tab:scaling}
\small
\begin{tabular}{lrrl}
\toprule
\textbf{Device class} & \textbf{Crypto (ms)} & \textbf{95th \%ile \Dt} &
  \textbf{Recommendation} \\
\midrule
Desktop, i7 AVX2   & $<$0.15   & $<$25    & Recommended \\
Raspberry Pi 4     & 3--10     & 20--35   & Use bounded pool \\
Low-power ARM SBC  & 12--25    & 35--60   & Conservative cap required \\
Cortex-M4 MCU      & 120--600  & $700+$   & Not recommended \\
\bottomrule
\end{tabular}
\end{table}

\subsection{Attack Vectors and Mitigations}

\paragraph{Classical MitM during \Dt.}
An adversary who can impersonate the server using a fraudulently obtained ECDSA
certificate passes ECDSA verification and obtains a provisional session; PQC
verification subsequently detects the invalid $\sigma_P$ and terminates the
connection.  Data exchanged during \Dt\ may be exposed.  The $D_\text{policy}$
cap directly bounds this exposure; applications should restrict provisional
transmission to non-sensitive content.

\paragraph{CRQC-enabled impersonation.}
An adversary with quantum capability to forge ECDSA could also forge a valid
PQC signature, defeating both verification stages.  This adversary is outside
the threat model of the provisional window; full quantum-safe authentication
holds only after \Dt\ elapses.

\paragraph{Handshake transcript replay.}
Reuse of a captured \texttt{CertificateVerify} is prevented by TLS~1.3's
session-specific transcript hash, which incorporates ephemeral key material.
The \texttt{Finished} MAC additionally binds $\sigma_P$ bytes before any
application data is transmitted.

\paragraph{Provisional flooding (DAPV-specific DoS).}
An adversary can flood connections with a valid $\sigma_E$ but deliberately
malformed $\sigma_P$.  Each connection enters provisional state, allocates a
background thread, and fails only when that thread executes.  With an unbounded
thread pool, this attack exhausts server resources.  A bounded pool with fixed
maximum worker count $M$, supplemented by connection-rate limiting, contains
this attack: the adversary can occupy at most $M$ background workers
simultaneously while legitimate foreground connections continue to be served.

\paragraph{Exfiltration via large provisional transfers.}
If $D_\text{policy}$ is misconfigured to allow arbitrarily large provisional
transmissions, an adversary in provisional state can exfiltrate substantial data.
The cap mitigates this directly; a value of $\leq 8$~KB is recommended based on
the 95th-percentile \Dt\ at 80\% CPU load.

Table~\ref{tab:mitigations} summarises the mitigation mapping.

\begin{table}[htbp]
\centering
\caption{DAPV attack vectors and mitigations.}
\label{tab:mitigations}
\small
\begin{tabular}{p{3.8cm}p{6.0cm}}
\toprule
\textbf{Attack vector} & \textbf{Mitigation} \\
\midrule
Classical MitM during \Dt
  & $D_\text{policy} \leq 8$\,KB; restrict provisional window to
    non-sensitive content. \\[4pt]
CRQC impersonation
  & Outside current scope; full quantum-safe authentication guaranteed
    post-\Dt. \\[4pt]
Transcript replay
  & Inherent in TLS~1.3 transcript hash and \texttt{Finished} MAC. \\[4pt]
Provisional flooding (DoS)
  & Bounded thread pool (fixed $M$); connection rate limiting. \\[4pt]
Exfiltration via large transfer
  & $D_\text{policy}$ cap enforced at application layer
    (Algorithm~\ref{alg:policy}). \\
\bottomrule
\end{tabular}
\end{table}

\subsection{Application-Layer Policy Interface}

Implementing $D_\text{policy}$ enforcement requires the TLS library to expose
the current session authentication state to the application layer.  The
following design achieves this within OpenSSL~3.x via the existing
\texttt{SSL\_CTX\_set\_info\_callback} mechanism, requiring no changes to TLS
record processing or key derivation:

\begin{verbatim}
#define SSL_DAPV_STATE_NONE        0
#define SSL_DAPV_STATE_PROVISIONAL 1
#define SSL_DAPV_STATE_FULL_AUTH   2
#define SSL_DAPV_STATE_PQC_FAILED  3

int SSL_get_dapv_state(const SSL *ssl);

void dapv_info_cb(const SSL *ssl, int where, int ret) {
    if (where & SSL_CB_DAPV_PROVISIONAL)
        app_set_state(SSL_get_app_data(ssl), PROVISIONAL);
    else if (where & SSL_CB_DAPV_FULL_AUTH)
        app_set_state(SSL_get_app_data(ssl), FULL_AUTH);
    else if (where & SSL_CB_DAPV_PQC_FAILED)
        app_set_state(SSL_get_app_data(ssl), FAILED);
}

int handle_request(SSL *ssl, Request *req) {
    if (SSL_get_dapv_state(ssl) == SSL_DAPV_STATE_PROVISIONAL
        && req->body_size > DPOLICY_CAP)
        return defer_until_full_auth(req);
    return process(req);
}
\end{verbatim}

The implementation requires two new callback flags and one four-line getter,
with an estimated footprint of approximately 200 lines across
\texttt{ssl/ssl\_lib.c} and \texttt{include/openssl/ssl.h}.

%% ============================================================
\section{Deployment Considerations}
\label{sec:deployment}

\subsection{Applicable Scenarios}

DAPV is most beneficial where two conditions hold: (i)~the application can
tolerate a bounded interval of classical-level authentication, and (ii)~handshake
latency materially affects user experience or system throughput.  Content
delivery networks, cloud API gateways, edge load balancers, and interactive web
services satisfy both conditions.  Backend services processing authentication
credentials, financial transactions, or personally identifiable information
should retain synchronous PQC verification.

\subsection{Risk-Stratified Deployment Policy}

A practical deployment partitions endpoints by data sensitivity:
\begin{itemize}[noitemsep]
\item \textit{Low-sensitivity} (static assets, read-only APIs):
  DAPV with $D_\text{policy} = 8$~KB.
\item \textit{Moderate-sensitivity} (interactive API responses):
  DAPV with $D_\text{policy} = 4$~KB; buffer credential-bearing operations.
\item \textit{High-sensitivity} (authentication, payments, PII):
  synchronous PQC verification; DAPV disabled.
\end{itemize}

\subsection{Operational Requirements}

Secure deployment requires four conditions to hold.
First, the TLS library must expose \textsc{provisional} and
\textsc{fully\_authenticated} states to the application via the interface
described in Section~\ref{sec:security}.
Second, the thread pool must be strictly bounded; unbounded creation enables
the provisional flooding attack.
Third, verification timing metrics (\Dt\ mean and 95th percentile, PQC failure
rate, provisional session count) must be instrumented continuously;
anomalies may indicate active exploitation.
Fourth, $D_\text{policy}$ must be calibrated against 95th-percentile \Dt\
under production load: Section~\ref{sec:res_delta_t} demonstrates that load
jitter can expand \Dt\ by up to 2.5$\times$ relative to idle conditions.

\subsection{Standardisation Pathway}

IETF standardisation of DAPV as a TLS~1.3 extension would require specifying:
(i)~a \texttt{ClientHello} extension advertising DAPV capability;
(ii)~a \texttt{ServerHello} acknowledgement enabling deferred verification;
(iii)~a post-handshake signal upon PQC verification completion; and
(iv)~transcript binding requirements and provisional data constraints.  The
compatibility of DAPV with TLS~1.3's existing transcript integrity mechanisms
positions it as a viable candidate for an IETF draft within the TLS working
group.

%% ============================================================
\section{Comparison with Prior Work}
\label{sec:comparison}

Table~\ref{tab:priorwork} situates the present results within the PQC TLS
performance literature.  The synchronous baseline measurements replicate prior
reported values closely, validating the experimental methodology.  DAPV reduces
the latency overhead from 52\% to 12\% at 150\ms\ RTT and the throughput loss
from 55\% to 20\%, representing the first architectural mitigation of comparable
scope evaluated in the literature.

\begin{table}[htbp]
\centering
\caption{Comparison with prior PQC TLS performance studies.  Latency overhead
  reported as percentage increase over classical ECDSA at the closest
  comparable RTT.  ``---'' indicates no mitigation was proposed.}
\label{tab:priorwork}
\small
\resizebox{\columnwidth}{!}{%
\begin{tabular}{lllrrr}
\toprule
\textbf{Study} & \textbf{Scheme} & \textbf{Platform} &
  \textbf{Latency overhead} & \textbf{Throughput loss} & \textbf{Mitigation} \\
\midrule
Paquin et al.~\cite{Paquin2020}
  & Dilithium hybrid & x86-64 & +80--300\% & $>$50\% & --- \\
Sikeridis et al.~\cite{Sikeridis2020a}
  & ML-DSA hybrid & x86-64 & +52\% @ 150\ms & 55\% & --- \\
Montenegro et al.~\cite{Montenegro2026}
  & Various hybrids & x86-64 & +45--280\% & 40--60\% & --- \\
Sosnowski et al.~\cite{Sosnowski2023}
  & PQC-only & x86-64 & +60--250\% & 35--55\% & --- \\
\textbf{This work (sync)} & P-256+ML-DSA-44 & x86-64 (Win) & +52\% @ 150\ms & 55\% & --- \\
\textbf{This work (DAPV)} & P-256+ML-DSA-44 & x86-64 (Win) & +12\% @ 150\ms & 20\% & DAPV \\
\bottomrule
\end{tabular}%
}
\end{table}

%% ============================================================
\section{Conclusion}
\label{sec:conclusion}

This paper has presented Deferred Asynchronous PQC Verification (DAPV), a
protocol-level optimisation for TLS~1.3 that decouples the two verification
stages of a hybrid handshake.  By completing the handshake provisionally on
the basis of a fast ECDSA check and deferring the PQC verification to a
background thread, DAPV removes post-quantum verification computation from the
handshake critical path.  Controlled experiments confirm a constant latency
reduction of approximately 20\ms\ across all tested RTT values (0--150\ms) and
a server throughput improvement of 30--36\% for hybrid P-256\,+\,ML-DSA-44,
both independent of network round-trip time.

Decomposition of the Window of Vulnerability reveals that the 14--21\ms\
provisional window in the prototype is dominated by thread scheduling overhead
rather than cryptographic computation, which costs less than 0.15\ms.  An
optimised bounded thread-pool implementation would reduce \Dt\ by approximately
two orders of magnitude, substantially narrowing the provisional security
exposure.  A novel denial-of-service vector specific to deferred verification
architectures is characterised and its mitigation through pool bounding and
connection rate limiting is specified.

The contributions of this work extend a growing body of evidence that the
principal obstacle to PQC deployment at Internet scale is the structural
mismatch between post-quantum signature sizes and network MTU constraints.
DAPV addresses the computational component of this challenge and provides a
deployable near-term optimisation for content-delivery, cloud, and edge
workloads while the fragmentation problem is pursued through complementary
network-layer and cryptographic-engineering approaches.

\subsection*{Future Work}

\begin{enumerate}[noitemsep]
\item \textit{Formal security proof}: game-based EU-CMA analysis of DAPV's
  authentication guarantees under combined adversarial models.
\item \textit{Optimised concurrency}: production-grade thread-pool and
  event-driven dispatcher implementations targeting \Dt\ below 1\ms.
\item \textit{Heterogeneous platforms}: evaluation on ARM Cortex-A, RISC-V,
  and hardware-accelerated PQC devices; adaptation for DTLS~1.3.
\item \textit{PQC-only provisional anchoring}: design of a lightweight anchor
  mechanism for deployments without a classical signature component.
\item \textit{IETF standardisation}: development of a TLS~1.3 extension draft
  specifying DAPV negotiation, state signalling, and transcript binding.
\item \textit{Extended algorithm evaluation}: SLH-DSA (SPHINCS+, FIPS~205) and
  MAYO and other NIST additional signature candidates.
\end{enumerate}

%% ---- Acknowledgements ----------------------------------------
\section*{Acknowledgements}

The authors thank the Open Quantum Safe project and the OpenSSL community for
the open-source infrastructure that made this work possible.

\section*{Declaration of Competing Interest}

The authors declare no competing financial interests or personal relationships
that could have influenced the work reported in this paper.

\bibliographystyle{elsarticle-num}
\bibliography{references}

\end{document}
